% !TEX program = xelatex
% PyQCU / QUDA MultiGrid semantic and formal benchmark report.
% The figures are intentionally generated by pgfplots from the recorded
% values, so the report has no dependency on a raster-image converter.
\documentclass[UTF8,a4paper,11pt]{ctexart}

\usepackage[a4paper,margin=2.05cm,headheight=14pt]{geometry}
\usepackage{amsmath,amssymb,mathtools,bm}
\usepackage{booktabs,longtable,array,tabularx,multirow}
\usepackage{graphicx}
\usepackage{xcolor}
\usepackage{listings}
\usepackage{url}
\usepackage{hyperref}
\usepackage{pgfplots}
\usepgfplotslibrary{groupplots}
\pgfplotsset{compat=1.18}

\definecolor{pyblue}{RGB}{36,91,163}
\definecolor{qudared}{RGB}{190,65,55}
\definecolor{stagegray}{RGB}{90,105,115}
\hypersetup{colorlinks=true,linkcolor=pyblue,urlcolor=pyblue,citecolor=pyblue,
  pdftitle={QUDA 与 PyQCU Strict MultiGrid 语义与正式对照}}
\urlstyle{tt}
\newcommand{\pathref}[1]{\path{#1}}
\newcommand{\code}[1]{\texttt{\detokenize{#1}}}
\newcommand{\dd}{\mathrm{d}}
\newcommand{\tr}{\operatorname{tr}}
\newcommand{\C}{\mathbb C}
\newcommand{\R}{\mathbb R}
\newcommand{\norm}[1]{\left\lVert#1\right\rVert}
\newcommand{\ip}[2]{\left\langle#1,#2\right\rangle}
\newcommand{\fine}{\mathrm f}
\newcommand{\coarse}{\mathrm c}
\newcommand{\pc}{\mathrm{PC}}
\newcommand{\odd}{\mathrm o}
\newcommand{\even}{\mathrm e}
\newcommand{\Yhat}{\widehat Y}
\newcommand{\Hhat}{\widehat H}
\newcommand{\Dhat}{\widehat D}
\newcommand{\todo}[1]{\textcolor{red}{\textbf{[待核验：#1]}}}
\setlength{\parindent}{2em}
\setlength{\parskip}{0.35em}
\emergencystretch=3em

\title{QUDA 与 PyQCU MultiGrid 的实现语义及正式大格对照}
\author{PyQCU 工程分析记录}
\date{2026 年 9 月 6 日}

\begin{document}
\maketitle
\begin{abstract}
本文以 \pathref{/root/PyQCU} 中的 PyQCU 和
\pathref{/root/PyQCU/refer/git-rep/quda} 中的 QUDA 1.1.0 源码为对象，
从 Wilson--Clover 奇偶块矩阵的第一性原理定义出发，逐层解释聚合型
MultiGrid 的 null vector、aggregate、延拓矩阵 $P$、限制矩阵 $R$、
Galerkin 粗算子、粗层 $X/Y/\Yhat$、Clover/Gauge、MATPC、平滑器、
粗层求解器和外层 Krylov 求解器。PyQCU 中同时存在旧的纯 Python/CUDA
33 点 Schur 路径、参考性的 compact Schur 层级和新的 Strict full-coarse
路径；三者在本文中严格分开。最后复现固定 V100、
$16\times32\times32\times48$、c64、odd--odd MATPC 协议，给出五次 steady
solve 的总时间、外层迭代数、逐迭代残差和 PyQCU 阶段计时，并把当前 QUDA
公开接口不能提供逐 GCR/V-cycle 阶段事件这一事实单独标出。
\end{abstract}

\tableofcontents

\section{范围、结论与证据等级}

本报告的工作对象是 PyQCU 的 MultiGrid 相关代码，主要对照
\pathref{/root/PyQCU/refer/git-rep/quda} 中的
\pathref{lib/transfer.cpp}、\pathref{lib/dirac_clover.cpp}、
\pathref{lib/dirac_coarse.cpp}、\pathref{lib/multigrid.cpp} 以及粗算子和
transfer CUDA kernel。结论分为三类：

\begin{itemize}
  \item ``源码确证''：可由下文给出的文件和行号直接定位；
  \item ``代数推导''：由已确证的矩阵块和存储约定推出，并给出推导链；
  \item ``未完成的逐元素核验''：不能因为正式求解收敛就冒充为
        $Y/\Yhat$ 每个方向、每个 parity、每个周期边界位置都已经和 QUDA
        逐元素相等。
\end{itemize}

核心结论可以压缩为如下数据流：
\begin{equation}
 (U_0,C_0)\;\xrightarrow{\;D_0=A_0-\kappa H_0\;}
 V_0\;\xrightarrow{\;R_0 A_0^{\rm eff}P_0\;}
 (X_1,Y_1,X_1^{-1},\Yhat_1)\;\xrightarrow{\;V_1\;}
 (X_2,Y_2,X_2^{-1},\Yhat_2)\;\cdots .
 \label{eq:data-flow}
\end{equation}
这里的 $C_0$ 是 fine Clover，$U_0$ 是 fine Gauge；level $1$ 以后使用的是
有效的矩阵 block，而不是把 SU(3) Gauge ``下采样''。一般 Galerkin 关系为
\begin{equation}
 D_{\ell+1}=R_\ell A_\ell^{\rm eff}P_\ell,
 \qquad R_\ell=P_\ell^\dagger,
 \qquad
 A_\ell^{\rm eff}\in\{D_\ell,\;X_\ell^{-1}D_\ell\}.
 \label{eq:galerkin-summary}
\end{equation}

本次正式结果的结论不是 ``每一次迭代 QUDA 都更快''，而是：在固定的
V100、编译变体、输入 bundle 和计时协议下，PyQCU Strict 的总 solve-only
中位数略快，但单次外层迭代更慢、外层迭代更少；该比值不具有跨平台普适性。

\section{统一记号、布局和量纲}

\subsection{格点、奇偶和自由度}

物理坐标按 $(x,y,z,t)$ 排列，所有时空轴在 PyQCU 约定中位于最后四维。
周期边界下
\begin{equation}
 p(x,y,z,t)=(x+y+z+t)\bmod 2,
 \qquad \even=0,\qquad \odd=1 .
 \label{eq:parity}
\end{equation}
一次 Wilson hopping 改变一个坐标奇偶，因此只连接 $\even\leftrightarrow\odd$；
Clover 是站点内的 spin--color 矩阵，不改变 parity。fine Wilson/Clover 的
每站点自由度为
\begin{equation}
 e=N_sN_c=4\times3=12.
 \label{eq:fine-dof}
\end{equation}
本次正式协议为 $n_v=12$、coarse spin $N_s^c=2$，所以每个 coarse site
拥有
\begin{equation}
 E=N_s^c n_v=2\times12=24
 \label{eq:coarse-dof}
\end{equation}
个粗自由度。这里 $n_v$ 是每一个 coarse-spin block 的 coarse color 数，
不是把 fine spin 直接删掉。

严格路径中，粗层资产使用完整几何
$[E,X_c,Y_c,Z_c,T_c]$；只有一个 parity 的场使用 compact 几何
$[E,X_c,Y_c,Z_c,T_c/2]$。在 QCU C++ 的 blocked transfer 布局中，
fine basis 的静态形状是
\begin{equation}
 [E,e,X_c,b_x,Y_c,b_y,Z_c,b_z,T_c,b_t],
 \qquad (X_f,Y_f,Z_f,T_f)=(X_cb_x,Y_cb_y,Z_cb_z,T_cb_t).
 \label{eq:blocked-layout}
\end{equation}
因此 parity 的 ``半格''不能被误读为粗资产本身只有半个格点。

\subsection{fine Gauge、Clover 和无量纲检查}

令 $U_\mu(x)$ 位于 $x\to x+\hat\mu$ 的有向链上，则未吸收系数的
Wilson hopping 可写为
\begin{equation}
 (H\psi)(x)=\sum_{\mu=0}^{3}\left[
 (1-\gamma_\mu)U_\mu(x)\psi(x+\hat\mu)
 +(1+\gamma_\mu)U_\mu^\dagger(x-\hat\mu)\psi(x-\hat\mu)
 \right].
 \label{eq:wilson-hopping}
\end{equation}
fine Dirac 算子采用
\begin{equation}
 D_0=A_0-\kappa H_0,
 \qquad
 A_0=\operatorname{diag}(A_{0,\even},A_{0,\odd}),
 \label{eq:fine-dirac}
\end{equation}
其中每个 $A_{0,p}(x)$ 是 $12\times12$ Clover onsite block。于是按
$(\even,\odd)$ 排列
\begin{equation}
 D_0=\begin{pmatrix}
 A_{0,\even}&-\kappa H_{\even\odd}\\
 -\kappa H_{\odd\even}&A_{0,\odd}
 \end{pmatrix} .
 \label{eq:block-fine}
\end{equation}
QUDA 的直接 full operator 在
\pathref{refer/git-rep/quda/lib/dirac_clover.cpp:35--65} 调用
\code{ApplyWilsonClover}，Clover-only 在 :50--56，C++ 严格路径则把
\code{clover_ee}、\code{clover_oo} 及两个 inverse 传进
\pathref{cpp/cuda/qcu/src/apply_multigrid_strict.cu:1818--1841}。

格点单位已把晶格间距吸收进参数，$U$、$A$、$H$、$D$、$X$、$Y$ 都是无量纲
矩阵；$P/R$ 只改变表示空间，不改变场的物理量纲。正式配置使用
\begin{equation}
 \kappa=\frac{1}{2(m+4)}=\frac{1}{2m+8},
 \qquad m=0.05,
 \qquad \frac{1}{2\kappa}=m+4=4.05 .
 \label{eq:kappa}
\end{equation}
最后一个等式也是本次 full residual 归一化中 PyQCU 与 QUDA solution scale
差异的锚点；它不是 Clover 方向错误。

\section{三条 PyQCU 路径与 QUDA 路径}

代码库中 ``MultiGrid'' 不是只有一个类。把下表的路径混为一谈，会直接把
33 点 Schur stencil、Strict 的最近邻 $X/Y$ 层级和 QUDA 的 native coarse
operator 当成同一个 ABI。

\begin{table}[htbp]
\centering
\caption{实现路径的边界和实际算子}
\small
\begin{tabularx}{\textwidth}{@{}p{0.18\textwidth}p{0.22\textwidth}X p{0.20\textwidth}@{}}
\toprule
路径 & 层级算子 & parity / 粗算子存储 & 主要求解器与源码 \\
\midrule
旧版 \code{multigrid} & fine parity 可选；粗层由
\code{dslash.operator} 和 33 点 stencil 组成 & 细层可先做 odd Schur；
粗层 stencil 为 onsite、8 个轴向、24 个双轴位移 & 每层 BiCGStab，
\code{cycle()} 在计数达到阈值时做粗修正；
\pathref{pyqcu/solver/_multigrid.py:353--540} \\
\midrule
\code{QudaMultigrid} legacy/compact & 支持 full-field Galerkin，也支持
\code{setup_operator="schur"} 的紧凑 parity Schur & compact Schur 的
粗自由度作为一个整体；小问题可 materialize dense，通用 coarse 有
\code{QudaCoarseOperator}/\code{QudaMatPCOperator} & 可切换 MR/FGMRES 等
参考实现；\pathref{pyqcu/solver/_quda_multigrid.py:909--1968,1971--2041} \\
\midrule
Strict & 每层先构造完整 coarse 的 $X/Y$；调用边界再做 parity crop &
full coarse geometry，运行期保留 $X^{-1}$ 和 forward/backward $\Yhat$；
Strict setup 对非最近邻支撑 fail-closed & C++ 持久层级、融合右预条件
FGMRES；\pathref{cpp/cuda/qcu/src/apply_multigrid_strict.cu:1340--2466} \\
\midrule
QUDA native & \code{DiracClover}/\code{DiracCloverPC} 递归到
\code{DiracCoarse}/\code{DiracCoarsePC} & Transfer full/coarse-site subset
由运行时选择；粗层 $Y/X/X^{-1}/\Yhat$ 是 native GaugeField & MG factory
组合 smoother、coarse solver、V-cycle/recursive；
\pathref{refer/git-rep/quda/lib/multigrid.cpp:273--430,564--640} \\
\bottomrule
\end{tabularx}
\end{table}

\subsection{旧版纯 Python/CUDA 33 点路径}

\pathref{pyqcu/solver/_multigrid.py:353--540} 的 \code{cycle()} 对当前层
做 BiCGStab；当 $\code{count_restart}>\code{num_restart}$ 且尚有下一层时，
依次执行
\begin{equation}
 r\;\xrightarrow{R_\ell}\;r_{\ell+1}
 \xrightarrow{\text{recursive solve}}e_{\ell+1}
 \xrightarrow{P_\ell}e_\ell,
 \qquad x_\ell\leftarrow x_\ell+e_\ell,
 \qquad r_\ell\leftarrow b_\ell-A_\ell x_\ell .
 \label{eq:legacy-cycle}
\end{equation}
如果 \code{support_parity=True}，level 0 先用
\code{give_b_parity} 形成 odd Schur RHS，求解后由
\code{give_x_e} 恢复 even 分量。C++ restrict/prolong 只在
\code{with_cuda_qcu} 时替换 Python einsum；粗 dslash 的 packing 为
\code{[2,4,E,E,X,Y,Z,T]}，见
\pathref{pyqcu/solver/_multigrid.py:91--199}。

\pathref{pyqcu/tools/_multigrid.py:277--875} 的 Schur 粗算子不是最近邻
9-block，而是
\begin{equation}
 A^{33}_c=S_c^{(0)}
 +\sum_{\mu=0}^{3}\sum_{\sigma=\pm1}S_c^{(\sigma\hat\mu)}
 +\sum_{0\leq\mu<\nu\leq3}\sum_{\sigma,\tau=\pm1}
 S_c^{(\sigma\hat\mu+\tau\hat\nu)} .
 \label{eq:33-stencil}
\end{equation}
计数为 $1+8+6\times4=33$。原因是 fine odd Schur 中存在两次 hopping：
两个相反方向给 onsite，两个同轴方向给 $\pm\hat\mu$，两个不同轴给
四种 $\pm\hat\mu\pm\hat\nu$。因此 33 点只适用于已经选定的 Schur
粗化路径；不能把它直接塞进 Strict 的仅支持 $X+\sum_{\pm\mu}Y_\mu$
ABI。代码中的 \code{PAIRS} 和 \code{SIGN} 正是 $6$ 对轴和每对 $4$
个符号组合。

多卡 \code{MultiGpuMultigrid} 在
\pathref{pyqcu/cuda/_multi_gpu.py:63--168} 为每个线程复制
\code{params/argv/set_ptrs}，在主线程或多线程构造 Schur 33 点资产并以
h5py 缓存；worker 线程主要执行各自卡上的 C++ 求解。它要求单 MPI rank，
并且当前 PyTorch wheel 对 P100 \code{sm_60} 的 torch kernel 有架构限制，
所以不能拿该路径的 worker 约束替代 Strict 的 full-coarse 语义。

\subsection{compact Schur 参考路径}

\code{CompactParityOperator} 将 full field 映射到一个 compact parity，默认
消去 even、暴露 odd Schur：
\begin{equation}
 S_\odd=D_{\odd\odd}-D_{\odd\even}A_{\even}^{-1}D_{\even\odd},
 \qquad
 b_\odd^S=b_\odd-D_{\odd\even}A_{\even}^{-1}b_\even .
 \label{eq:compact-schur}
\end{equation}
它在 \pathref{pyqcu/solver/_quda_multigrid.py:1762--1881} 中显式实现
\code{apply/rhs/reconstruct}；\code{QudaMatPCOperator} 则保留完整 coarse
资产，在应用时通过 checkerboard embed/extract 做两次跨 parity hopping，
见 \pathref{pyqcu/solver/_quda_multigrid.py:1539--1595}。这条路径适合小规模
代数等价性测试和参考矩阵构造；它和 Strict 的区别是：Strict 的
\code{X/Y/Yhat} setup 必须先在 full coarse geometry 上完成，compact parity
只在 solver 边界出现。

\subsection{Strict full-coarse 路径}

\code{QudaMultigrid} 在 \code{hierarchy_mode="strict"} 时要求
\code{use_parity=True}，并在 \pathref{pyqcu/solver/_quda_multigrid.py:2089--2227}
拒绝当前未实现的 staggered/KD、奇数 coarse extent 和不一致的 smoother
模式。\code{CudaStrictMultigridSolver} 在
\pathref{pyqcu/cuda/_strict_multigrid.py:90--226} 中完成 MPI preflight、
setup、资产绑定、持久 hierarchy 初始化和显存预算；其运行入口在
\pathref{pyqcu/cuda/_strict_multigrid.py:267--334}。

Strict 的重要工程不变量为：
\begin{equation}
 \code{applyInitQcu}
 \longrightarrow \text{绑定 level assets}
 \longrightarrow \text{操作}
 \longrightarrow \code{params[_SET_INDEX_] += 1}
 \longrightarrow \code{applyEndQcu} .
 \label{eq:qcu-lifetime}
\end{equation}
每个实例自己的 \code{params/argv/set_ptrs} 必须独立；否则 scratch slot
和 CUDA stream 之间会发生缓冲复用冲突。Strict 的 stage trace 是可选的，
默认关闭，因此不改变无 trace 生产路径的算子顺序。

\section{从 fine Schur 到粗层 $P/R$ 和 Galerkin}

\subsection{奇偶 Schur 的完整推导}

从式 \eqref{eq:block-fine} 的第二行得到
\begin{equation}
 x_q=A_q^{-1}\left(b_q+\kappa H_{qp}x_p\right),
 \qquad q=1-p .
 \label{eq:reconstruct-fine}
\end{equation}
代入第一行得到非对称 Schur 方程
\begin{equation}
 \left(A_p-\kappa^2H_{pq}A_q^{-1}H_{qp}\right)x_p
 =b_p+\kappa H_{pq}A_q^{-1}b_q .
 \label{eq:asym-schur}
\end{equation}
若左乘 $A_p^{-1}$，则得到对称归一化 MATPC 算子
\begin{equation}
 M_p=A_p^{-1}\left(A_p-\kappa^2H_{pq}A_q^{-1}H_{qp}\right)
 =I-\kappa^2A_p^{-1}H_{pq}A_q^{-1}H_{qp},
 \label{eq:sym-schur}
\end{equation}
以及其 RHS
\begin{equation}
 b_p^{\pc}=A_p^{-1}\left(b_p+\kappa H_{pq}A_q^{-1}b_q\right).
 \label{eq:prepare-fine}
\end{equation}
这里 $p$ 可以是 even 或 odd；本次 formal 协议取 $p=\odd$。QUDA
\code{DiracCloverPC::M} 在
\pathref{refer/git-rep/quda/lib/dirac_clover.cpp:173--210} 分开处理
asymmetric 和 symmetric 分支；\code{prepare/reconstruct} 在
\pathref{refer/git-rep/quda/lib/dirac_clover.cpp:223--260} 直接对应
式 \eqref{eq:prepare-fine} 和 \eqref{eq:reconstruct-fine}。

Strict C++ 的链条是
\begin{align}
 H_{qp}x_p&\xrightarrow{A_q^{-1}}A_q^{-1}H_{qp}x_p
 \xrightarrow{H_{pq}}H_{pq}A_q^{-1}H_{qp}x_p
 \xrightarrow{A_p^{-1}}A_p^{-1}H_{pq}A_q^{-1}H_{qp}x_p, \\
 M_px_p&=x_p-\kappa^2A_p^{-1}H_{pq}A_q^{-1}H_{qp}x_p,
 \end{align}
对应 \pathref{cpp/cuda/qcu/src/apply_multigrid_strict.cu:1843--1863}。
如果把目标 parity 的 $A_p^{-1}$ 提前到 RHS 括号内部之前，得到的将是
$A_p^{-1}b_p+\kappa H_{pq}A_q^{-1}b_q$，并非式
\eqref{eq:prepare-fine}；该乘法顺序是必须守住的语义边界。

\subsection{aggregate、null vector 和 block Gram--Schmidt}

给定 block $b=(b_x,b_y,b_z,b_t)$，fine 坐标映射到
\begin{equation}
 X_\mu=\left\lfloor x_\mu/b_\mu\right\rfloor,
 \qquad
 \Lambda_{\ell+1}=\left(L_x/b_x,L_y/b_y,L_z/b_z,L_t/b_t\right).
 \label{eq:aggregate-map}
\end{equation}
Wilson/Clover 的 fine spin $4\to2$ 映射为
\begin{equation}
 s_c=\left\lfloor s_f/2\right\rfloor,
 \qquad s_f=0,1\mapsto s_c=0,quad s_f=2,3\mapsto s_c=1 .
 \label{eq:spin-map}
\end{equation}
每个 coarse site 和 coarse-spin block 内，对输入的 near-null vectors 做
局部 Euclidean block CGS。若 $B_j$ 是第 $j$ 列，重复两遍时
\begin{equation}
 v_j^{(r)}=B_j^{(r)}-sum_{i<j}v_i^{(r)}
 \ip{v_i^{(r)}}{B_j^{(r)}}_{\mathcal A},
 \qquad
 v_j^{(r)}\leftarrow
 \frac{v_j^{(r)}}{\sqrt{\ip{v_j^{(r)}}{v_j^{(r)}}_{\mathcal A}}},
 \qquad r=1,2 .
 \label{eq:block-cgs}
\end{equation}
在 PyQCU 当前参考代码中 $\ip{\cdot}{\cdot}_{\mathcal A}$ 实现为该
aggregate 内的标准复 Euclidean 内积；符号只强调 ``局部块内''，不是
另行引入一个未在代码中构造的 Clover 加权内积。实现为
\pathref{pyqcu/solver/_quda_multigrid.py:316--346,446--482}，QUDA 的
对应入口为 \code{Transfer::reset -> BlockOrthogonalize}，见
\pathref{refer/git-rep/quda/lib/transfer.cpp:152--163}。

把正交后的基记为
$V_{s_f,c_f,s_c,j}(x)$。对 coarse field $\phi$ 的延拓和对 fine field
$\psi$ 的限制为
\begin{align}
 (P\phi)_{s_f,c_f}(x)&=
 \sum_{j=0}^{n_v-1}V_{s_f,c_f,s_c,j}(x)\,
 \phi_{s_c,j}(X(x)), \\
 (R\psi)_{s_c,j}(X)&=
 \sum_{x\in\mathcal A_X}\sum_{s_f,c_f}
 V_{s_f,c_f,s_c,j}(x)^*\psi_{s_f,c_f}(x), \\
 R&=P^\dagger .
 \label{eq:PR}
\end{align}
PyQCU 的 full-field 实现是
\pathref{pyqcu/solver/_quda_multigrid.py:538--597,610--658}，
parity view 则在 :673--696；QUDA prolongator 先做 geo/spin map 再做
fine-color rotation，见
\pathref{refer/git-rep/quda/include/kernels/prolongator.cuh:61--127}，
restrictor 使用共轭基，见
\pathref{refer/git-rep/quda/include/kernels/restrictor.cuh:81--124}。

最重要的两个投影检查是
\begin{equation}
 \ip{P\phi}{\psi}=\ip{\phi}{R\psi},
 \qquad
 P^\dagger P=I_{\text{coarse subspace}},
 \qquad
 PP^\dagger\ne I_{\text{fine space}}\quad\text{一般成立}.
 \label{eq:projection-check}
\end{equation}
第二个等式不是说 fine 空间被完整反演，而是说每个局部正交 coarse basis
在其张成子空间内归一。

\subsection{Galerkin block 的构造}

对任意有效 fine 算子 $A_\ell^{\rm eff}$，定义
\begin{equation}
 D_{\ell+1}=R_\ell A_\ell^{\rm eff}P_\ell .
 \label{eq:galerkin}
\end{equation}
Strict direct-PC 选择
\begin{equation}
 A_\ell^{\rm eff}=X_\ell^{-1}D_\ell,
 \qquad
 D_{\ell+1}=R_\ell(X_\ell^{-1}D_\ell)P_\ell .
 \label{eq:strict-galerkin}
\end{equation}
PyQCU 在 \pathref{pyqcu/solver/_quda_multigrid.py:2695--2815} 根据
\code{hierarchy_mode}、\code{smoother_solve_type} 和
\code{strict_galerkin_mode} 选择这个算子；batched strict builder 在
\pathref{pyqcu/tools/_strict_galerkin.py:594--720} 对每一批 source aggregate
置入 $E$ 个 blocked basis，调用 batch matvec，再计算 $V^\dagger A V$。

固定 source aggregate $X$，只有 source block 和一个 fine nearest-neighbor
所能到达的相邻 aggregate 会非零。令 $\Delta$ 属于这些 displacement，则
\begin{equation}
 B_\Delta(X)=\left(P_\ell^\dagger A_\ell^{\rm eff}P_\ell\right)
 _{X,X+\Delta},
 \qquad
 X_{\ell+1}=B_0,
 \qquad
 Y_{\ell+1}^{\pm\hat\mu}=B_{\pm\hat\mu} .
 \label{eq:block-projection}
\end{equation}
Strict builder 的 support check 要求图像在 source aggregate 和八个 one-hop
target aggregate 外为零，见
\pathref{pyqcu/tools/_strict_galerkin.py:678--699}。这是 Strict 最近邻
X/Y ABI 的适用条件，也是 33 点 Schur 路径不能直接复用的原因。

\section{粗层 $X/Y/\Yhat$、Clover-like dslash 与矩阵顺序}

\subsection{有效粗矩阵的 storage 约定}

在目标点 action-coefficient 记号下，可以写成
\begin{equation}
 (D_c z)(X)=X(X)z(X)+\sum_\mu\left[
 Y_\mu^{\rm f}(X)z(X+\hat\mu)
 +Y_\mu^{\rm b}(X)z(X-\hat\mu)\right].
 \label{eq:coarse-action}
\end{equation}
QUDA 的 native storage 使用 link-like 约定：forward $Y^{\rm f}$ 放在源点，
backward link 放在 $X-\hat\mu$，gather 时取 dagger。对应的 CUDA 注释是
\pathref{refer/git-rep/quda/include/kernels/dslash_coarse.cuh:126--141}，
backward bulk gather 在 :232--249，onsite coarse Clover block 在
:257--290。因而不能只对 tensor 做一个 `roll` 就假定方向已经正确，
必须同时确定 storage site、gather site、dagger 和矩阵左右乘顺序。

对 coarse-PC，令
\begin{equation}
 \Dhat_c=X_c^{-1}D_c=I+\Hhat_c .
 \label{eq:coarse-pc}
\end{equation}
需要保存
\begin{equation}
 \begin{split}
 \Yhat_\mu^{\rm f}(X)&=X_c^{-1}(X)Y_\mu^{\rm f}(X),\\
 \Yhat_\mu^{\rm b}(X-\hat\mu)&=
 Y_\mu^{\rm b}(X-\hat\mu)X_c^{-\dagger}(X-\hat\mu).
 \end{split}
 \label{eq:yhat}
\end{equation}
QUDA 的逐 kernel 证据在
\pathref{refer/git-rep/quda/include/kernels/coarse_op_preconditioned.cuh:48--121}：
backward 分支先做 $Y X^{-\dagger}$，forward 分支做 $X^{-1}Y$；
\code{DiracCoarsePC::createCoarseOp} 在
\pathref{refer/git-rep/quda/lib/dirac_coarse.cpp:628--647} 明确说递归粗化
的是 $\Yhat$ 而不是 $Y$。

PyQCU 的 Strict pack 在
\pathref{pyqcu/tools/_strict_galerkin.py:489--525} 先批量求 $X^{-1}$，
将 backward action coefficient 转到 $q-\hat\mu$ storage，再分别计算
\code{X_inv @ forward} 和 \code{backward_storage @ X_inv^†}；
\code{QudaCoarseOperator} 的 runtime 绑定在
\pathref{pyqcu/solver/_quda_multigrid.py:1265--1305}。因此 $\Yhat$ 不是对
已经完成的粗 dslash 输出再 ``乘一次逆''；误写为
$X^{-1}Y^{\rm b}$ 会改变非平凡 Clover 下的算子。

\subsection{粗 dslash、粗 Clover 和层间 Gauge 的消失}

粗层 dslash 读取的是 $E\times E$ 的有效 link-like block：
\begin{equation}
 (D_c z)(X)=X(X)z(X)+\sum_\mu\left[
 Y_\mu^{\rm f}(X)z(X+\hat\mu)
 +Y_\mu^{\rm b\dagger}(X-\hat\mu)z(X-\hat\mu)
 \right],
 \label{eq:coarse-native-dslash}
\end{equation}
其中把 backward dagger 写在 action 中只是为了与 QUDA 的 storage 对齐。
\code{DiracCoarse::Dslash} 由 \pathref{refer/git-rep/quda/include/kernels/dslash_coarse.cuh:140--255}
执行 hop，\code{applyClover} 执行 onsite block；
\code{DiracCoarsePC::Dslash} 在
\pathref{refer/git-rep/quda/lib/dirac_coarse.cpp:506--519} 只把输入 link
换成 $\Yhat$，onsite $X$ 仍显式参与。

因此层间对象的类型变化是
\begin{equation}
 (U_0\in SU(3),C_0\in\C^{12\times12})
 \longrightarrow
 (X_1,Y_1\in\C^{24\times24})
 \longrightarrow
 (X_2,Y_2\in\C^{E_2\times E_2}) .
 \label{eq:no-gauge-coarse}
\end{equation}
原始 Gauge/Clover 只定义 fine operator 和第一个投影；后续层的 ``Clover''
是 $X_\ell$，``Gauge'' 是 $Y_\ell$ 或 $\Yhat_\ell$ 的命名类比，而不是
SU(3) 矩阵。QUDA 的 \code{DiracCloverPC::createCoarseOp} 入口见
\pathref{refer/git-rep/quda/lib/dirac_clover.cpp:263--270}，粗层再由
\code{DiracCoarse/DiracCoarsePC} 接管。

\section{奇偶边界：MATPC、prepare 和 reconstruct}

\subsection{full geometry 与 compact parity 的关系}

每个 coarse level 的资产保留
\begin{equation}
 \mathcal G_\ell=\{X_\ell,Y_\ell^{\rm f},Y_\ell^{\rm b},
 X_\ell^{-1},\Yhat_\ell^{\rm f},\Yhat_\ell^{\rm b}\}
 \quad\text{on full }\Lambda_\ell .
 \label{eq:full-assets}
\end{equation}
compact parity 只定义为 full field 的一个视图。Strict C++ 的
\code{strict_restrict_parity_kernel} 与 \code{strict_prolong_parity_kernel}
接收完整 coarse $X,Y,Z,T$，在
\pathref{cpp/cuda/qcu/src/apply_multigrid_strict.cu:700--820} 处理 parity。
若固定前三维，compact 时间索引可以表示为
\begin{equation}
 t_{\rm full}=2t_c+\left[(p-x-y-z)\bmod2\right] .
 \label{eq:compact-time}
\end{equation}
这解释了为什么 ``数组第二列''、``物理 odd site'' 和 ``checkerboard index''
不是同一个概念。当前 Strict 还要求每个非平凡 coarse extent 为偶数，
否则周期 checkerboard 没有等量的两种 parity。

\subsection{coarse MATPC 的两次 hopping}

用式 \eqref{eq:coarse-pc} 的归一化表示，目标 parity $p$ 的 coarse MATPC 为
\begin{equation}
 M_p^c=I-\Hhat_{pq}\Hhat_{qp},
 \qquad q=1-p .
 \label{eq:coarse-matpc}
\end{equation}
严格的应用顺序是
\begin{align}
 u_q&=\Hhat_{qp}x_p,\\
 v_p&=\Hhat_{pq}u_q,\\
 M_p^cx_p&=x_p-v_p .
 \label{eq:matpc-sequence}
\end{align}
Strict C++ 在 \pathref{cpp/cuda/qcu/src/apply_multigrid_strict.cu:2227--2241}
发射两次 \code{strict_hopping_parity_kernel}；PyQCU 参考类在
\pathref{pyqcu/tools/_strict_galerkin.py:378--404} 先 embed、两次 hopping、
再 extract。QUDA 的 odd--odd 分支在
\pathref{refer/git-rep/quda/lib/dirac_coarse.cpp:545--557}，其 full field
会直接报错，因为 coarse-PC operator 要求 parity subset。

\subsection{fine prepare/reconstruct 与 coarse prepare/reconstruct}

对 fine 的符号约定 $D=A-\kappa H$，prepare/reconstruct 已由式
\eqref{eq:prepare-fine} 和式 \eqref{eq:reconstruct-fine} 给出。对粗层，
QUDA 把 off-diagonal block 的符号吸收到粗 $Y$ 或 kernel 的 $\kappa$ 参数中，
所以源码写作
\begin{equation}
 b_p^{c,\pc}=X_p^{-1}b_p^{c}
 -\left(X_p^{-1}D_{pq}^{c}\right)X_q^{-1}b_q^{c},
 \qquad
 x_q^c=X_q^{-1}b_q^c-
 \left(X_q^{-1}D_{qp}^{c}\right)x_p^c .
 \label{eq:coarse-prepare}
\end{equation}
若把粗 off-diagonal 定义成 $-D_{pq}^c$，式中两个减号会变成加号；
代数对象不变，storage 符号必须和 kernel 同时读取。QUDA 的实际分支为
\pathref{refer/git-rep/quda/lib/dirac_coarse.cpp:570--625}。

\section{平滑器、粗求解器和外层预条件}

\subsection{MR 平滑器}

Strict 的 fine 和 coarse MR 使用同一个非 Hermitian-safe 标量更新。给定
当前 residual $r$，令 $v=Mr$，其中 $M$ 是当前 parity operator，则
\begin{equation}
 \alpha=\frac{\ip{v}{r}}{\ip{v}{v}},
 \qquad x\leftarrow x+\alpha r,
 \qquad r\leftarrow r-\alpha v .
 \label{eq:mr}
\end{equation}
对应 \pathref{cpp/cuda/qcu/src/apply_multigrid_strict.cu:1928--1946}
和 :2283--2300。分母 finite/breakdown guard 不是装饰：coarse Schur 或
方向相关 $\Yhat$ 不必在普通 Euclidean 内积下 Hermitian positive definite，
因此不能不加条件地替换成 CG。

QUDA 的 smoother factory 在
\pathref{refer/git-rep/quda/lib/multigrid.cpp:273--332} 设置
\code{inv_type}、tol、Schwarz cycle、CA basis 以及是否返回 residual；
它还分别建立 residual operator、smoother operator 和 sloppy smoother。
当前 formal 配置使用 $\nu_{pre}=\nu_{post}=1$，但代码枚举可表达的 MR、
CG、CA-CG、GCR、BiCGStab(l) 或 Schwarz 不能被说成在本次正式 run 中全部
执行过。

\subsection{coarsest BiCGStab}

Strict 最粗层使用 BiCGStab；以下表格按用户要求把整轮递推写成单列伪代码，
并明确 preconditioned/operator 的位置。令 $A=M_L$ 为最粗层 MATPC，
本表中 $M=1$ 表示不再嵌套第二个 preconditioner。

\begin{table}[htbp]
\centering
\caption*{算法 1：最粗层 BiCGStab（带可选左/右预处理的单列伪代码）}
\small
\begin{tabular}{@{}l@{}}
$r^{(0)}=b-Ax^{(0)}$，$\widetilde r=r^{(0)}$，$p$、$v$、$s$、$t$ 置零。\\
for $i=1,2,\ldots$，直到 $\norm{r^{(i)}}^2\le\tau^2\norm b^2$：\\
\quad$\rho_{i-1}=\widetilde r^\dagger r^{(i-1)}$；若 $|\rho_{i-1}|$ 或 $|\omega_{i-1}|$ 太小，则报告 breakdown。\\
\quad 若 $i=1$，$p^{(1)}=r^{(0)}$；否则
$\beta_{i-1}=\dfrac{\rho_{i-1}}{\rho_{i-2}}
\dfrac{\alpha_{i-1}}{\omega_{i-1}}$，
$p^{(i)}=r^{(i-1)}+\beta_{i-1}(p^{(i-1)}-\omega_{i-1}v^{(i-1)})$。\\
\quad solve $M\widehat p=p^{(i)}$；若无预处理，取 $M=I$、$\widehat p=p^{(i)}$。\\
\quad$v^{(i)}=A\widehat p$，
$\alpha_i=\dfrac{\rho_{i-1}}{\widetilde r^\dagger v^{(i)}}$。\\
\quad$s=r^{(i-1)}-\alpha_i v^{(i)}$；若 $\norm{s}\le\tau\norm b$，
$x^{(i)}=x^{(i-1)}+\alpha_i\widehat p$，然后停止。\\
\quad solve $M\widehat s=s$；$t=A\widehat s$；
$\omega_i=\dfrac{t^\dagger s}{t^\dagger t}$。\\
\quad$r^{(i)}=s-\omega_i t$，
$x^{(i)}=x^{(i-1)}+\alpha_i\widehat p+\omega_i\widehat s$。\\
\quad 若 $\rho$、分母或 $\omega$ 非 finite，则保留当前 finite iterate，
重置 shadow/search/scalar 状态或按调用方策略退出。\\
repeat until convergence or $i=i_{\max}$；记录 $i$、$\norm{r^{(i)}}/\norm b$ 和耗时。\\
\end{tabular}
\end{table}

Strict C++ 的实际最粗层对应
\pathref{cpp/cuda/qcu/src/apply_multigrid_strict.cu:2303--2378}；它采用
relative squared stopping target、finite guard 和 $\rho/\alpha/\omega$ 复数
递推。legacy Python BiCGStab 还在粗修正后重置整个递推状态，见
\pathref{pyqcu/solver/_multigrid.py:487--501}；这是因为 $r$ 已被粗修正改变。

\subsection{V-cycle、W-cycle、F-cycle 与 recursive}

一次 V-cycle 的误差传播可写成
\begin{equation}
 \mathcal V_\ell(b)=S_{\ell,post}^{\nu_{post}}
 \left[\,S_{\ell,pre}^{\nu_{pre}}(0,b)
 +P_\ell C_{\ell+1}R_\ell
 \left(b-A_\ell S_{\ell,pre}^{\nu_{pre}}(0,b)\right)\right].
 \label{eq:v-cycle}
\end{equation}
W-cycle 在同一 child 上调用两次，F-cycle 的第一次下行更强，recursive
则把 child MG object 直接交给 coarse solver wrapper。QUDA 当前
\code{createCoarseSolver} 的可执行 factory 分支主要是 VCYCLE 或 RECURSIVE，
证据为 \pathref{refer/git-rep/quda/lib/multigrid.cpp:564--573}；枚举中存在
F/W 并不等价于本次 formal 已经运行过 F/W。

Strict 的两层实际执行顺序是：
\begin{equation}
 \text{fine MR}_{pre}\to R\to
 \text{level-1 prepare}\to\text{level-1 BiCGStab}
 \to\text{reconstruct}\to P\to\text{fine MR}_{post} .
 \label{eq:strict-v-sequence}
\end{equation}
C++ 的递归实现 \pathref{cpp/cuda/qcu/src/apply_multigrid_strict.cu:2407--2466}
还会在 child 返回后重新计算 parent residual，而不是保存每一层三份 parent
大小的向量。

\subsection{右预条件 FGMRES}

Strict formal 外层求解 fine odd Schur，预条件器是整个 MR--V-cycle--MR：
\begin{equation}
 r_0=b-M_px_0,
 \qquad z_j=\mathcal M_j^{-1}v_j,
 \qquad w_j=M_pz_j,
 \label{eq:right-fgmres}
\end{equation}
其中 $\mathcal M_j^{-1}$ 可因层级、精度、guard 和内部状态而改变，故使用
flexible 而不是假定固定线性 $M^{-1}$ 的普通 GMRES。Arnoldi/Givens 更新为
\begin{equation}
 h_{ij}=\ip{v_i}{w_j},
 \qquad
 w_j\leftarrow w_j-\sum_{i=0}^{j}h_{ij}v_i,
 \qquad
 v_{j+1}=w_j/\norm{w_j},
 \qquad
 x\leftarrow x+\sum_j z_j y_j .
 \label{eq:fgmres-recurrence}
\end{equation}
Strict C++ 在 \pathref{cpp/cuda/qcu/src/apply_multigrid_strict.cu:2070--2215}
实现该递推，并在 restart boundary 必做 true residual refresh；正式 run 的
requested restart 为 16，但 512 MiB workspace 约束把 effective restart
缩到 4。

\section{一段贯穿全流程的超长公式型伪代码}

下面的单列长表把从输入 Gauge/Clover、null vector 到 full solution 的因果
链写在同一处。它同时标出 ``直接 full Galerkin''、``Schur/33 点'' 和
``Strict/nearest-neighbor'' 的分叉；这样从零复现时不会在中途把某一条路径
的 parity 或 storage 约定替换到另一条路径。

\begin{longtable}{@{}p{0.965\textwidth}@{}}
\caption{算法 2：PyQCU/QUDA MultiGrid 的 setup--apply--solve 全流程单列伪代码}\label{tab:long-algorithm}\\
\toprule
输入：周期 fine lattice $\Lambda_0=(L_x,L_y,L_z,L_t)$、Gauge $U_0$、Clover
$A_{0,\even},A_{0,\odd}$、$\kappa$、target parity $p$、block $b_\ell$、
每层 null vector 数 $n_{v,\ell}$、coarse spin、smoother、coarse solver、
层数 $L$、outer tolerance $\tau$。\\
按式 \eqref{eq:parity} 给每个 site 标记 parity；确认所有用于 checkerboard
的非平凡 extent 为偶数，并区分 full index、parity index、compact index。\\
构造 $D_0=A_0-\kappa H(U_0)$；Clover 只在 site 内作用，Gauge 只在式
\eqref{eq:wilson-hopping} 的 forward/backward hopping 中作用。\\
若配置 full solve，设 $q=1-p$；若配置 MATPC solve，输入输出只保存 $p$
compact field，但仍保留 $q$ 的 RHS、Clover inverse 和 hopping workspace。\\
若 null vector 已由 canonical bundle/cache 提供，读取并验证 shape、dtype、
layout、content digest；否则生成随机/test/inverse/Krylov 起点并以 setup
solver 反复平滑，使 $B\approx D^{-1}$ 的低模方向。\\
对 $\ell=0,1,\ldots,L-2$：按 $X_\mu=\lfloor x_\mu/b_{\ell,\mu}\rfloor$
建立 fine-to-coarse 与 coarse-to-fine map，令 $\Lambda_{\ell+1}$ 为 block 后
的完整几何。\\
对 Wilson/Clover 设置 $s_c=\lfloor s_f/2\rfloor$；对每个 aggregate、
每个 coarse-spin block，执行重复 block CGS 和归一化，得到
$V_\ell(s_f,c_f,s_c,j,x)$。\\
定义 $P_\ell$ 为 $V_\ell$ 的 coarse-to-fine contraction，定义
$R_\ell=P_\ell^\dagger$ 为 fine-to-coarse contraction；检查
$\ip{P\phi}{\psi}=\ip{\phi}{R\psi}$，但不要把 $PP^\dagger$ 当作 fine identity。\\
若路径是旧版 Schur：先在 fine parity 上形成 $S_\ell$，再令
$A_\ell^{\rm eff}=S_\ell$；由于两次 hopping，准备 33 个 displacement
$(0,\pm\hat\mu,\pm\hat\mu\pm\hat\nu)$。\\
若路径是 full/Strict：令 $A_\ell^{\rm eff}=D_\ell$ 或
$X_\ell^{-1}D_\ell$；在 Strict + direct-PC 中必须使用后者，不能先投影
$D_\ell$ 再事后任意补一个 $X_\ell^{-1}$。\\
对每个 coarse source $X$ 探测 $P_\ell^\dagger A_\ell^{\rm eff}P_\ell$；
把零 displacement 记为 $X_{\ell+1}$，把 $+\hat\mu/-\hat\mu$ 记为
$Y_{\ell+1}^{\rm f/b}$，把 Schur 双轴 displacement 记入 33 点 stencil。\\
若为 Strict，检查 probe image 在 source aggregate 和八个 one-hop target
aggregate 外为零；若出现双轴或更宽支撑，拒绝把该算子写入 Strict X/Y ABI。\\
逐 coarse site 批量求 $X_{\ell+1}^{-1}$；形成
$\Yhat^{\rm f}=X^{-1}Y^{\rm f}$，形成
$\Yhat^{\rm b}=Y^{\rm b}X^{-\dagger}$，并把 backward link 放到
$X-\hat\mu$ storage；执行 periodic wrap 和 halo exchange。\\
丢弃或缓存 raw $Y$ 取决于路径：Strict solve 可只保留 $X^{-1}$、$\Yhat$、
blocked $V$；诊断和后续 coarse construction 可以额外保留 raw $Y$。\\
若 $\ell+1$ 仍需继续粗化，下一层输入不再是原始 $U_0/C_0$，而是本层有效
$(X,Y)$ 或 coarse-PC 的 $(X,\Yhat)$；更新 coarse spin、$n_v$、dof 和几何。\\
setup 结束后，构造 residual operator、smoother operator、sloppy operator、
presmoother、postsmoother 和 coarse solver；QUDA 中这些对象可以是不同精度
或不同 solve type，不能用一个 ``operator'' 名称代替全部。\\
给定 full RHS $b=(b_p,b_q)$ 和初值 $x_0$，先执行
$b_p^{\pc}=A_p^{-1}(b_p+\kappa H_{pq}A_q^{-1}b_q)$；若只要 MATPC 解则
停止在 $x_p$，若要 full 解则保留 reconstruct 所需的 $b_q$。\\
设置 $r_0=b_p^{\pc}-M_px_{p,0}$；若 $\norm{r_0}$ 已满足容差，跳过 Krylov
但仍按 full-solution contract reconstruct 并计算 full true residual。\\
对 outer FGMRES 的每个 restart：令 $\beta=\norm{r}$、$v_0=r/\beta$，
清零 Hessenberg/Givens/scalar，并保存 flexible preconditioned basis $z_j$。\\
对 $j=0,1,\ldots$：先做 $\nu_{pre}$ 次 fine MR；把当前 residual 通过
$R_0$ 限制到完整 coarse geometry；调用 child V-cycle/recursive coarse solve。\\
child solver 内部对非最粗层执行 prepare、pre-MR、restriction、child solve、
prolongation、parent residual recompute、post-MR；最粗层执行表
\textnormal{算法 1} 的 BiCGStab 或配置的 GCR/CA-GCR。\\
将 coarse correction 通过 $P_0$ 延拓回 fine target parity，做 fine correction
residual，再做 $\nu_{post}$ 次 MR，得到 $z_j=\mathcal M_j^{-1}v_j$。\\
用 $w_j=M_pz_j$，计算 $h_{ij}=v_i^\dagger w_j$，正交化、归一化并应用
complex Givens；把当前 least-squares estimate 与 basis floor 写入 trace。\\
若 estimate 满足阈值或发生 happy breakdown，退出 inner cycle；回代小型
上三角系统得到 $y_j$，并用 $x_p\leftarrow x_p+\sum_jz_jy_j$ 更新解。\\
在每个 restart boundary 都重新算 $r=b_p^{\pc}-M_px_p$；Strict 中该
true-residual refresh 是强制的，不能只相信递推 residual。\\
若 $\norm r>\tau\norm b$ 且 outer iteration 未达到上限，继续下一个 restart；
若收敛，进入 reconstruct。\\
reconstruct：按 $x_q=A_q^{-1}(b_q+\kappa H_{qp}x_p)$ 恢复 fine 被消去
parity；粗层按其 storage 符号执行对应 $X_q^{-1}$ 和 $Y_q$ 组合。\\
构造 full $x=(x_p,x_q)$，以同一份 periodic Wilson/Clover full operator 计算
$\eta_{\rm true}=\norm{b-D_0x}_2/\norm b_2$；只以它作为跨实现的正确性 gate。\\
保存每次 outer iteration 的 residual kind、wall time、stage time、level、
restart boundary、最终 full true residual 和配置 hash；将 timing 和 trace
分开，避免 trace 同步开销污染正式 solve-only 中位数。\\
\bottomrule
\end{longtable}

\section{QUDA 原生实现的源码级流程}

\subsection{Transfer 与粗化 factory}

QUDA 的 \code{Transfer} 构造中先校验 aggregate size 与 $Nvec$，再分配
\code{V}、fine/coarse map 和 spin map，见
\pathref{refer/git-rep/quda/lib/transfer.cpp:23--114}；\code{createV}
在 :117--149 把 coarse vector 数编码到 native field order，\code{reset}
在 :152--163 调用 block orthogonalization。\code{createGeoMap} 在
:200--230 把每个 fine site 映射到 parity-ordered coarse site。

\code{MG::createCoarseDirac} 选择用于构造粗算子的
\code{diracResidual} 或 \code{diracSmoother}：
\begin{equation}
 \text{preconditioned\_coarsen}
 \Longleftrightarrow
 (\text{coarse solution}=\text{MATPC})
 \land(\text{smoother solve}=\text{DIRECT\_PC}) .
 \label{eq:quda-coarsen-choice}
\end{equation}
源码在 \pathref{refer/git-rep/quda/lib/multigrid.cpp:345--421}，并根据之前
层是否已经 preconditioned 决定 \code{need_bidirectional}。这说明 QUDA
不是单纯 ``每层重做相同的 dslash''：residual/smoother/PC 的选择会改变
粗化输入和必须保留的方向信息。

\subsection{null-space setup 与递归}

没有从文件加载时，\code{MG::generateNullVectors} 在
\pathref{refer/git-rep/quda/lib/multigrid.cpp:1275--1320} 设置 setup solver、
tol、precision、CA basis 和 \code{compute_null_vector}，然后用 solver
更新 null vector；粗层 null vector 也可由 $R_\ell B_\ell$ 后重新正交得到。
本次 formal 对照固定使用已经转换并验证的 12 个 canonical full vectors，
所以不能把 formal setup 时间解释为 ``两侧都在线生成了同样的 null vectors''。

\code{MG::operator()} 首先按 outer MATPC type 选择 parity，并按
\code{coarse_grid_solution_type} 切换 full/parity site subset，证据为
\pathref{refer/git-rep/quda/lib/multigrid.cpp:1131--1160}。若 smoother 和
coarse correction 的 solution type 不一致，QUDA 会显式重算 residual；
这正是 residual operator 与 smoother operator 必须分开的工程原因。

\section{PyQCU Strict 的逐层执行与资产生命周期}

\subsection{setup 的批量构造}

Strict setup 的主要步骤是：
\begin{enumerate}
  \item \code{QudaTransfer} 生成 full-field $V$，并导出 QCU blocked layout；
  \item \code{build_strict_galerkin} 以 colored/site batch 方式计算
        $R(X^{-1}D)P$，检查 nearest-neighbor support；
  \item \code{_finish_strict_galerkin} 生成 $(X,X^{-1})$、raw $Y$（可选）
        和 forward/backward $\Yhat$；
  \item \code{qcu_strict_transition_assets} 把各 transition 的四槽资产
        写进 \code{set_ptrs}；
  \item \code{applyMultigridStrictInitQcu} 创建持久 hierarchy，并根据
        \code{(2m+5)B_f+2B_c} 守住 outer workspace 上限。
\end{enumerate}

形式上，严格 builder 的 workspace 由一个粗点批次和 $E$ 个 colored probe
共同决定；正式数据记录了 coarse sites $=49{,}152$、coarse dof $=24$、
colored column batch $=12$、projection site batch $=4$、operator calls $=44$。
这些统计在 \pathref{data/strict_vs_quda_formal_20260906.json}，不是从旧
Markdown 复制的估计值。

\subsection{Strict C++ 的一次外层迭代}

\code{precondition_fine} 的实际阶段在
\pathref{cpp/cuda/qcu/src/apply_multigrid_strict.cu:1948--1994}：
\begin{equation}
 \begin{array}{c}
 v\xrightarrow{\text{fine MR}_{pre}}u
 \xrightarrow{R_0}b_1
 \xrightarrow{\text{coarse solve}}e_1
 \xrightarrow{P_0}u+P_0e_1\\[-2pt]
 \xrightarrow{M_p}\text{correction residual}
 \xrightarrow{\text{fine MR}_{post}}z .
 \end{array}
 \label{eq:strict-preconditioner-flow}
\end{equation}
\code{solve_level} 在 :2407--2466 处理子层；
\code{apply_matpc/prepare/reconstruct} 在 :2227--2281 只通过完整 geometry
的 assets 做 parity kernel。外层 \code{fgmres_impl} 每个 restart 强制执行
\code{true_residual_refresh}，并将 trace stage 围在 CUDA stream sync 之间；
因此 stage 时间包含诊断同步开销，仅用于归因。

\subsection{Strict 与 33 点的不可替换边界}

Strict 的输入 operator 是 full nearest-neighbor $D$ 或左预条件后的
$X^{-1}D$，粗 block 支持为 zero plus eight axis neighbors；旧版 Schur
路径的输入是 parity Schur，天然包含二次 hopping 和 33 个位移。若对
Schur 只保存 Strict 的 $X/Y$，双轴贡献会被静默丢弃，所得对象即使在平凡
Gauge 上偶尔收敛，也不再是原 Galerkin 算子。这是 Strict builder 进行
support check 而不是盲目截断的数学理由。

\section{正式大格复现协议与结果}

\subsection{协议}

所有 timing 使用无 trace benchmark 的 steady samples；trace 仅收集
逐外层 residual 和 PyQCU stage。正式配置如下。

\begin{table}[htbp]
\centering
\caption{正式 benchmark 固定协议}
\small
\begin{tabularx}{\textwidth}{@{}p{0.28\textwidth}X@{}}
\toprule
项目 & 固定值 \\
\midrule
硬件 & Tesla V100-SXM2-32GB，compute capability 7.0，单卡单 MPI rank；
UUID 前缀 \code{be23...f0a8} \\
软件 & PyTorch \code{2.10.0+cu128}，CUDA 12.8，PyQUDA 0.10.54；
QUDA 1.1.0 WSL2 \code{DEV87_REDUCE_SYNC} 构建 \\
格点与物理 & $(L_x,L_y,L_z,L_t)=(16,32,32,48)$，$m=0.05$，
$\kappa=0.1234567901234568$，seed 42，周期边界 \\
层级 & 2 levels，block $(2,2,2,2)$，$n_v=12$，coarse spin $=2$，
coarse dof $=24$ \\
奇偶 & target parity $p=1$，QUDA \code{QUDA_MATPC_ODD_ODD} \\
外层 & restarted right FGMRES/GCR；requested restart $=16$，effective restart
$=4$，max outer iter $=1000$，tol $=10^{-6}$ \\
平滑/粗求解 & $\nu_{pre}=\nu_{post}=1$；coarse maxiter $=200$，
coarse tol $=3\times10^{-3}$ \\
统计 & 2 次 warmup + 5 次 steady；报告 median 和 MAD，warmup 不计入 \\
正确性 & full periodic Wilson/Clover true residual gate
$\eta_{true}<5\times10^{-6}$ \\
身份校验 & \code{config_hash=654b7267...332478}，
\code{bundle_hash=8c866bd0...a6ca3d31} \\
\bottomrule
\end{tabularx}
\end{table}

workspace 量纲检查为
\begin{align}
 B_f&=12\times16\times32\times32\times24\times8
   =37{,}748{,}736\ \mathrm{B},\\
 B_c&=24\times8\times16\times16\times24\times8
   =9{,}437{,}184\ \mathrm{B},\\
 B_{\rm Krylov}&=(2\times4+5)B_f+2B_c
   =509{,}607{,}936\ \mathrm{B}<512\ \mathrm{MiB} .
 \label{eq:workspace}
\end{align}

\subsection{solve-only 总结果}

\begin{table}[htbp]
\centering
\caption{五次 steady solve：外层总迭代、solve-only 时间和 full true residual}
\small
\begin{tabular}{@{}lrrrrrrr@{}}
\toprule
侧 & solve 0 & solve 1 & solve 2 & solve 3 & solve 4 & median (s) & MAD (s) \\
\midrule
PyQCU Strict & $11/2.021366$ & $11/2.110346$ & $11/2.015191$ &
$11/2.041653$ & $11/2.068720$ & $2.041653$ & $0.026462$ \\
QUDA & $37/2.094076$ & $37/2.094841$ & $37/2.082558$ &
$37/2.123275$ & $37/2.108544$ & $2.094841$ & $0.012283$ \\
\bottomrule
\end{tabular}
\end{table}
表内每格是 ``外层迭代数 / 秒数''。于是
\begin{equation}
 \frac{t_{\rm QUDA}}{t_{\rm PyQCU}}=1.026051,
 \qquad
 \frac{11}{37}=0.2973,
 \qquad
 \frac{185.605}{56.617}=3.278 .
 \label{eq:timing-ratios}
\end{equation}
PyQCU 的 outer iteration 较少，但单次 outer iteration 成本约为 QUDA 的
3.28 倍，所以总时间只快约 $2.6\%$。这只适用于本节的固定环境。

\begin{table}[htbp]
\centering
\caption{正确性、setup 和显存归因}
\small
\begin{tabularx}{\textwidth}{@{}p{0.29\textwidth}rrX@{}}
\toprule
指标 & PyQCU Strict & QUDA & 说明 \\
\midrule
5 次 steady full true residual & $3.6013\times10^{-7}$ &
$7.3030\times10^{-7}$ & 都低于 $5\times10^{-6}$ gate，5/5 converged \\
setup seconds & 14.777423 & 364.369869 & PyQCU 为 schema-v2 runtime cache hit；
QUDA 建立 native MG runtime \\
cache restore & 13.030576 & -- & PyQCU 其余 setup non-cache 为 1.746847 s \\
owned/runtime assets & 4,076,863,488 B & 未由 PyQUDA 暴露 & PyQCU 约 3.797 GiB \\
fused workspace & 509,607,936 B & native 未暴露 & PyQCU 约 0.475 GiB \\
setup device-wide peak & 11,219,046,400 B & 21,014,843,392 B & 设备级采样可能含 native/其他进程 \\
first-solve peak & 11,722,362,880 B & 24,529,670,144 B & lazy workspace；不计入 steady timing \\
\bottomrule
\end{tabularx}
\end{table}

\subsection{逐外层迭代 residual}

下表采用首个 steady solve 的完整曲线；五次 steady 的迭代数一致，完整
250 行逐迭代记录在 \pathref{data/strict_trace_detailed_20260906.csv}。
PyQCU 是 Arnoldi least-squares estimate，QUDA 是 GCR iterated residual，
两列不是同一个递推标量；最终判断仍使用 full true residual。

\begin{longtable}{@{}rll@{}}
\caption{首个 steady solve 的每次外层迭代 residual}\label{tab:residual-trace}\\
\toprule
$k$ & PyQCU Strict & QUDA \\
\midrule
\endfirsthead
\toprule
$k$ & PyQCU Strict & QUDA \\
\midrule
\endhead
0 & $1.0000000\times10^{0}$ & $1.0000000\times10^{0}$ \\
1 & $9.3650706\times10^{-2}$ & $2.6259840\times10^{-1}$ \\
2 & $1.6313400\times10^{-2}$ & $1.0671660\times10^{-1}$ \\
3 & $2.8122286\times10^{-3}$ & $5.6699110\times10^{-2}$ \\
4 & $6.4573408\times10^{-4}$ & $3.2547150\times10^{-2}$ \\
5 & $2.1028987\times10^{-4}$ & $2.1907580\times10^{-2}$ \\
6 & $6.2526138\times10^{-5}$ & $1.2670970\times10^{-2}$ \\
7 & $2.1975775\times10^{-5}$ & $8.4522200\times10^{-3}$ \\
8 & $7.4830368\times10^{-6}$ & $5.3343160\times10^{-3}$ \\
9 & $3.2432135\times10^{-6}$ & $3.8731680\times10^{-3}$ \\
10 & $1.2322430\times10^{-6}$ & $2.4618190\times10^{-3}$ \\
11 & $4.0650434\times10^{-7}$ & $1.7779130\times10^{-3}$ \\
12 & -- & $1.1777290\times10^{-3}$ \\
13 & -- & $8.7806670\times10^{-4}$ \\
14 & -- & $6.0349930\times10^{-4}$ \\
15 & -- & $4.3366300\times10^{-4}$ \\
16 & -- & $3.1065080\times10^{-4}$ \\
17 & -- & $2.3763540\times10^{-4}$ \\
18 & -- & $1.6512580\times10^{-4}$ \\
19 & -- & $1.2454290\times10^{-4}$ \\
20 & -- & $9.0326790\times10^{-5}$ \\
21 & -- & $6.8630000\times10^{-5}$ \\
22 & -- & $5.0507290\times10^{-5}$ \\
23 & -- & $3.7877360\times10^{-5}$ \\
24 & -- & $2.8157010\times10^{-5}$ \\
25 & -- & $2.2275490\times10^{-5}$ \\
26 & -- & $1.6195150\times10^{-5}$ \\
27 & -- & $1.2465930\times10^{-5}$ \\
28 & -- & $9.3341650\times10^{-6}$ \\
29 & -- & $7.4462140\times10^{-6}$ \\
30 & -- & $5.4595650\times10^{-6}$ \\
31 & -- & $4.2656100\times10^{-6}$ \\
32 & -- & $3.2030910\times10^{-6}$ \\
33 & -- & $2.5857410\times10^{-6}$ \\
34 & -- & $1.8816930\times10^{-6}$ \\
35 & -- & $1.5172770\times10^{-6}$ \\
36 & -- & $1.1199800\times10^{-6}$ \\
37 & -- & $9.1188520\times10^{-7}$ \\
\bottomrule
\end{longtable}

PyQCU 在第 11 轮的内部 estimate 为 $4.0650\times10^{-7}$；QUDA 在第
37 轮的 GCR iterated residual 为 $9.1189\times10^{-7}$。不能对这两列逐点
相减后宣称 ``PyQCU 误差''；可比较的是下降趋势和独立 full-op gate。

\subsection{阶段计时：哪些成本主导一次 Strict outer iteration}

PyQCU trace 共含 5 个 steady solve、每个 11 个 outer iteration；因为阶段
trace 会在 CUDA stream 前后同步，表中数值含同步开销，并且 nested stage
不能相加。QUDA 当前 \code{invertQuda/TimeProfile} 公开路径只给 aggregate
profile，没有逐 GCR 或逐 V-cycle 的回调，因此下表不伪造 QUDA 阶段值。

\begin{table}[htbp]
\centering
\caption{PyQCU Strict 阶段中位数（55 个 outer iteration 的阶段）}
\small
\begin{tabular}{@{}llrr@{}}
\toprule
level & stage & count & median (ms) \\
\midrule
0 & outer iteration & 55 & 185.852 \\
1 & coarse V-cycle & 55 & 164.769 \\
1 & coarsest BiCGStab & 55 & 157.980 \\
1 & level prepare & 55 & 3.247 \\
1 & level reconstruct & 55 & 2.923 \\
0 & fine pre-smoother & 55 & 3.257 \\
0 & fine restriction & 55 & 3.307 \\
0 & fine prolongation & 55 & 2.726 \\
0 & fine correction residual & 55 & 2.644 \\
0 & fine post-smoother & 55 & 3.104 \\
0 & fine MATPC (nested) & 55 & 2.476 \\
0 & Arnoldi (nested) & 55 & 1.319 \\
0 & FGMRES solution update & 15 & 0.630 \\
0 & true residual refresh & 15 & 3.090 \\
\bottomrule
\end{tabular}
\end{table}

其中 coarse V-cycle 占 outer iteration 的中位数约 $88.9\%$，coarsest
BiCGStab 本身约 $85.1\%$。这给出了比 ``总时间/总迭代'' 更有用的热点
判断：若要优化当前 formal，首先应分析最粗层递推、粗层矩阵应用和同步，
而不是只优化几毫秒的 $P/R$ kernel。

\begin{figure}[htbp]
\centering
\begin{tikzpicture}
\begin{axis}[
  width=0.88\textwidth,height=0.46\textwidth,
  xlabel={外层迭代 $k$}, ylabel={reported relative residual},
  xmode=linear, ymode=log, xmin=0,xmax=37,ymin=5e-7,ymax=2,
  grid=both, major grid style={gray!25}, minor grid style={gray!12},
  legend style={at={(0.98,0.98)},anchor=north east,font=\small},
  tick label style={font=\small}, label style={font=\small}]
\addplot[very thick,pyblue,mark=*,mark size=1.3pt] coordinates {
(0,1) (1,0.0936507061) (2,0.0163134001) (3,0.00281222863)
(4,0.000645734079) (5,0.000210289872) (6,0.0000625261382)
(7,0.0000219757749) (8,0.00000748303682) (9,0.00000324321354)
(10,0.00000123224299) (11,0.000000406504341)};
\addlegendentry{PyQCU Strict（Arnoldi estimate）}
\addplot[very thick,qudared,mark=square*,mark size=1.3pt] coordinates {
(0,1) (1,0.2625984) (2,0.1067166) (3,0.05669911)
(4,0.03254715) (5,0.02190758) (6,0.01267097) (7,0.00845222)
(8,0.005334316) (9,0.003873168) (10,0.002461819) (11,0.001777913)
(12,0.001177729) (13,0.0008780667) (14,0.0006034993) (15,0.000433663)
(16,0.0003106508) (17,0.0002376354) (18,0.0001651258) (19,0.0001245429)
(20,0.00009032679) (21,0.00006863) (22,0.00005050729) (23,0.00003787736)
(24,0.00002815701) (25,0.00002227549) (26,0.00001619515) (27,0.00001246593)
(28,0.000009334165) (29,0.000007446214) (30,0.000005459565)
(31,0.00000426561) (32,0.000003203091) (33,0.000002585741)
(34,0.000001881693) (35,0.000001517277) (36,0.00000111998) (37,0.0000009118852)};
\addlegendentry{QUDA（GCR iterated residual）}
\end{axis}
\end{tikzpicture}
\caption{首个 steady solve 的逐外层 residual。两条曲线的 scalar contract 不同，
图只用于展示各自递推的收敛行为。}
\label{fig:residual}
\end{figure}

\begin{figure}[htbp]
\centering
\begin{tikzpicture}
\begin{groupplot}[group style={group size=2 by 1,horizontal sep=1.1cm},
  width=0.42\textwidth,height=0.34\textwidth,
  tick label style={font=\small}, label style={font=\small}]
\nextgroupplot[ybar,bar width=18pt,ymin=0,ymax=2.35,
  ylabel={solve-only median (s)}, symbolic x coords={PyQCU,QUDA},
  xtick=data, nodes near coords, nodes near coords align={vertical},
  every node near coord/.append style={font=\scriptsize},
  title={总 solve 时间}]
\addplot[fill=pyblue!75,draw=pyblue] coordinates {(PyQCU,2.041653) (QUDA,2.094841)};
\nextgroupplot[ybar,bar width=18pt,ymin=0,ymax=210,
  ylabel={平均每外层迭代 (ms)}, symbolic x coords={PyQCU,QUDA},
  xtick=data, nodes near coords, nodes near coords align={vertical},
  every node near coord/.append style={font=\scriptsize},
  title={单次 outer iteration}]
\addplot[fill=qudared!75,draw=qudared] coordinates {(PyQCU,185.605) (QUDA,56.617)};
\end{groupplot}
\end{tikzpicture}
\caption{总时间和单次外层迭代时间必须同时看；只看后者会错误地忽略迭代数差异。}
\label{fig:time}
\end{figure}

\begin{figure}[htbp]
\centering
\begin{tikzpicture}
\begin{axis}[
  xbar, width=0.88\textwidth,height=0.52\textwidth,
  xlabel={阶段中位数 (ms)},
  symbolic y coords={true-res,FGMRES-update,Arnoldi,fine-MATPC,fine-post,
    fine-correction,fine-prolong,fine-restrict,fine-pre,level-reconstruct,
    level-prepare,coarsest-BiCGStab,coarse-V-cycle},
  ytick=data, xmin=0,xmax=175, enlarge y limits=0.04,
  nodes near coords, every node near coord/.append style={font=\scriptsize},
  tick label style={font=\scriptsize}, label style={font=\small},
  xmajorgrids, major grid style={gray!25}]
\addplot[fill=stagegray!75,draw=stagegray] coordinates {
  (3.090,true-res) (0.630,FGMRES-update) (1.319,Arnoldi)
  (2.476,fine-MATPC) (3.104,fine-post) (2.644,fine-correction)
  (2.726,fine-prolong) (3.307,fine-restrict) (3.257,fine-pre)
  (2.923,level-reconstruct) (3.247,level-prepare)
  (157.980,coarsest-BiCGStab) (164.769,coarse-V-cycle)};
\end{axis}
\end{tikzpicture}
\caption{PyQCU Strict 阶段诊断。嵌套 stage 不可相加；coarse-V-cycle 内含
coarsest-BiCGStab，且所有 stage 均含 trace 同步开销。}
\label{fig:stage}
\end{figure}

\section{Smoke、验证边界和未完成项}

\subsection{已完成的快速回归}

本轮已经运行并通过：
\begin{table}[htbp]
\centering
\caption{语义与实现快速闸门}
\small
\begin{tabular}{@{}lrrl@{}}
\toprule
类别 & 数量 & 结果 & 覆盖内容 \\
\midrule
CPU & 19 & PASS & Transfer、Galerkin、MATPC、prepare/reconstruct、FGMRES 代数 \\
CUDA Strict & 10 & PASS & full coarse、$X/Y/\Yhat$、parity kernel、生命周期 \\
融合 FGMRES & 3 & PASS & workspace budget、外层收敛、true residual \\
合计 & 32 & PASS & 无 skip/失败 \\
\bottomrule
\end{tabular}
\end{table}
结果摘要为 \pathref{data/strict_fast_cpu_20260906.json}；源代码语法检查和
build/install 回归也在本轮完成。测试不等价于证明 QUDA 与 PyQCU 的每一个
native storage element 相同，但证明了当前 Strict 的闭环可以运行并满足
固定协议的 full residual gate。

\subsection{QUDA reduction/MultiGrid smoke 的架构阻塞}

本轮另外执行了：
\begin{lstlisting}[basicstyle=\small\ttfamily,breaklines=true,frame=single]
python -B examples/qcu/dev87/smoke_quda_reduction.py \
  --quda-install data/quda-qio-install
python -B examples/qcu/dev87/smoke_quda_multigrid.py \
  --quda-install data/quda-qio-install
\end{lstlisting}
两条命令都在 QUDA 第一个 \code{CopyGauge} kernel 处报
\begin{lstlisting}[basicstyle=\small\ttfamily,breaklines=true,frame=single]
no kernel image is available for execution on the device
QUDA 1.1.0 (git 1.1.0--sm_60)
Using device 0: Tesla V100-SXM2-32GB
\end{lstlisting}
这表示该 smoke 所选的 \code{libquda.so} 是 \code{sm_60} 架构构建，而实际
设备为 V100 \code{sm_70}，且失败发生在进入 MultiGrid setup 之前。因此：
\begin{enumerate}
  \item smoke 结果应记录为 ``环境架构阻塞''，不能解释为算法失败；
  \item 不能把 smoke 的失败和已经记录的 formal JSON 结果拼成一个 ``当前
        所有 QUDA 路径都通过'' 的结论；
  \item 若要重新跑 smoke，应显式选择含 \code{sm_70} SASS/PTX 的 QUDA 安装，
        并重新记录库 hash、CMake features 和设备架构。
\end{enumerate}

formal JSON 的 QUDA 侧记录了另一份已完成 run 的 provenance，包括
\code{patch_variant=dev87_wsl2_reduce_sync}；它和当前 smoke 的架构选择是
两个必须分开的实验事实。

\subsection{当前未完成的逐元素核验}

仍应优先完成如下实验，而不应从本报告的 formal convergence 过度推断：
\begin{enumerate}
  \item 固定同一 canonical nontrivial Gauge/Clover 和同一 $V$，导出每个
        coarse site、每个方向、两个 parity 的 $X^{-1}$、$Y^{\rm f/b}$、
        $\Yhat^{\rm f/b}$；
  \item 对 backward storage 检查 $q-\hat\mu$ 的位置、gather dagger、
        periodic wrap 和 halo boundary；
  \item 用同一向量分别计算 QUDA native coarse dslash、PyQCU reference
        coarse dslash 和 Strict C++ kernel 的逐元素差；
  \item 同时比较 $P^\dagger P$、$RDP$、MATPC 两次 hopping 和 full
        reconstruct，而不仅比较最终 solver residual。
\end{enumerate}

另外，trace 残差不能替代 true residual；setup hit 的时间不能直接和 online
setup miss 比较；设备级 memory peak 可能包含 QUDA native allocation 或其他
进程，必须和 owned asset bytes 分开报告。

\section{可复现命令、产物和源码索引}

\subsection{命令}

\begin{lstlisting}[basicstyle=\small\ttfamily,breaklines=true,frame=single]
source ./env.sh
source examples/qcu/dev87/quda_env.sh

# formal solve-only timing（无 trace，2 warmup + 5 steady）
python -B examples/qcu/dev87/bench_strict_vs_quda.py \
  --profile formal --side both --cache-expect hit \
  --quda-nullvec-prefix data/L16x32x32x48_nvec12_quda \
  --quda-nullvec-manifest data/L16x32x32x48_nvec12_quda.conversion.json \
  --output data/strict_vs_quda_formal_20260906.json

# diagnostic trace（逐外层 residual 和 PyQCU stage）
python -B examples/qcu/dev87/trace_strict_vs_quda.py \
  --output data/strict_trace_stage_20260906.json \
  --plot data/strict_trace_stage_20260906.svg \
  --benchmark-output data/strict_trace_benchmark_20260906.json \
  --reference data/strict_vs_quda_formal_20260906.json \
  --repeats 5 --timeout 1800

# 展开 JSON，生成逐迭代/阶段 CSV、SVG 和摘要
python -B examples/qcu/dev87/analyze_strict_vs_quda_detailed.py
\end{lstlisting}

\subsection{本报告直接引用的产物}

\begin{itemize}
  \item formal：\pathref{data/strict_vs_quda_formal_20260906.json}；
  \item trace：\pathref{data/strict_trace_stage_20260906.json}；
  \item detailed：\pathref{data/strict_trace_detailed_20260906.json}、
        \pathref{data/strict_trace_detailed_20260906.csv}；
  \item stage：\pathref{data/strict_trace_stage_detailed_20260906.json}、
        \pathref{data/strict_trace_stage_timing_20260906.csv}；
  \item 自动图：\pathref{data/strict_trace_detailed_20260906.svg}；
  \item 快速回归：\pathref{data/strict_fast_cpu_20260906.json}；
  \item 分析器：\pathref{examples/qcu/dev87/analyze_strict_vs_quda_detailed.py}。
\end{itemize}

\subsection{关键源码索引}

\begin{longtable}{@{}p{0.28\textwidth}p{0.31\textwidth}p{0.31\textwidth}@{}}
\caption{从数学对象回到源码的索引}\label{tab:source-index}\\
\toprule
对象 & PyQCU & QUDA \\
\midrule
fine $D=A-\kappa H$ / Clover &
\pathref{cpp/cuda/qcu/src/apply_multigrid_strict.cu:1818--1900} &
\pathref{refer/git-rep/quda/lib/dirac_clover.cpp:35--65} \\
$P/R$、spin map、CGS &
\pathref{pyqcu/solver/_quda_multigrid.py:349--597,610--696} &
\pathref{refer/git-rep/quda/lib/transfer.cpp:23--163}；
\pathref{refer/git-rep/quda/include/kernels/prolongator.cuh:61--127} \\
full Galerkin / $X/Y$ &
\pathref{pyqcu/tools/_strict_galerkin.py:594--720}；
\pathref{pyqcu/solver/_quda_multigrid.py:909--1086} &
\pathref{refer/git-rep/quda/lib/coarse_op.in.cpp:8--20} \\
$X^{-1}$ / $\Yhat$ &
\pathref{pyqcu/tools/_strict_galerkin.py:489--525} &
\pathref{refer/git-rep/quda/include/kernels/coarse_op_preconditioned.cuh:48--121} \\
coarse dslash/Clover &
\pathref{pyqcu/solver/_quda_multigrid.py:1426--1501} &
\pathref{refer/git-rep/quda/lib/dirac_coarse.cpp:506--519}；
\pathref{refer/git-rep/quda/include/kernels/dslash_coarse.cuh:126--290} \\
MATPC / compact parity &
\pathref{pyqcu/solver/_quda_multigrid.py:1539--1689,1762--1881} &
\pathref{refer/git-rep/quda/lib/dirac_clover.cpp:173--260}；
\pathref{refer/git-rep/quda/lib/dirac_coarse.cpp:530--625} \\
legacy 33 点 &
\pathref{pyqcu/tools/_multigrid.py:277--875}；
\pathref{pyqcu/solver/_multigrid.py:353--540} & 不属于 QUDA native Strict X/Y ABI；
对应 native coarse 支持见 dslash/coarse-op kernel \\
Strict recursion/FGMRES &
\pathref{cpp/cuda/qcu/src/apply_multigrid_strict.cu:1948--2215,2227--2466} &
\pathref{refer/git-rep/quda/lib/multigrid.cpp:273--430,564--640,1131--1160} \\
\bottomrule
\end{longtable}

\section{总结：从零复现时应守住的六条不变量}

\begin{enumerate}
  \item 先按式 \eqref{eq:block-fine} 写清楚 $A$、$H$、$\kappa$ 和目标
        parity，再讨论 ``symmetric/asymmetric''；
  \item $R=P^\dagger$ 是 transfer 的代数不变量，$P/R$ 的局部 block map
        和 spin map 必须与字段 layout 同时验证；
  \item coarse $X/Y$ 是 Galerkin 矩阵 block，coarse ``Clover/Gauge'' 不是
        原始 $12\times12$ Clover 和 SU(3) Gauge 的下采样；
  \item backward $Y$ 的 storage site、dagger 和 $X^{-\dagger}$ 右乘顺序
        必须同时正确；
  \item full coarse assets 与 compact parity view 分离，MATPC 只在应用、
        prepare、reconstruct 边界 crop；
  \item solver 内部 residual、Arnoldi/GCR reported residual 和 full-op true
        residual 是三个对象，跨实现正确性只能以最后一个为准。
\end{enumerate}

本报告的正式数据、阶段数据、smoke 限制和未决逐元素核验均已落在工作区；
后续若优化 \code{strict_hopping_parity_kernel} 或改变粗层 storage，应重新
生成新的 config hash、formal trace 和 full true residual，而不能用本次结果
覆盖新语义。

\end{document}
